\documentclass{beamer}
\usepackage[T1]{fontenc}
\usepackage[utf8]{inputenc}
\usepackage{lmodern}
\usepackage{amsmath,amssymb,mathtools,bm}
\usepackage{booktabs}
\usepackage{listings}
\lstset{
    language=Matlab,
    basicstyle=\ttfamily\scriptsize,
    breaklines=true,
    columns=fullflexible,
    showstringspaces=false,
    frame=single
}
\usetheme{Boadilla}
\useoutertheme{miniframes}
\setbeamertemplate{navigation symbols}{}

\newcommand{\E}{\mathbb{E}}
\newcommand{\Q}{\mathbb{Q}}
\newcommand{\norm}[1]{\left\|#1\right\|}
\newcommand{\dd}{\mathrm{d}}
\newcommand{\wrong}[1]{\textcolor{red}{#1}}

\title{MATH4907 Midterm Test Review}
\author{Sanyang Liu}
\date{November 2026}

\AtBeginSection[]
{
    \begin{frame}{Table of Contents}
        \tableofcontents[currentsection,hideallsubsections]
    \end{frame}
}

\begin{document}
\maketitle

\section{Question 1: Concepts}

\begin{frame}{Q1(1): Arbitrage}
\small
\textbf{Select all correct answers.} Which statements about arbitrage
and the no-arbitrage principle are correct?
\begin{enumerate}
\item[A.] An arbitrage strategy may have non-positive initial value.
\item[B.] An arbitrage strategy has zero probability of loss at some future time.
\item[\wrong{C.}] \wrong{An arbitrage strategy must guarantee strictly positive profit in every state.}
\item[D.] If two portfolios have the same terminal payoff almost surely,
then they should have the same initial value under no-arbitrage.
\item[E.] No-arbitrage is one of the fundamental principles used to price
contingent claims.
\end{enumerate}
\end{frame}

\begin{frame}{Q1(1): Why?}
\small
\begin{itemize}
\item \textbf{A. True.} A zero-cost strategy can be an arbitrage;
a negative initial cost may also be possible.
\item \textbf{B. True.} An arbitrage has no chance of a terminal loss.
\item \wrong{\textbf{C. False.} Positive profit is required with positive
probability, not in every state.}
\item \textbf{D. True.} Different initial prices for the same payoff
would create a buy-low, sell-high arbitrage.
\item \textbf{E. True.} No-arbitrage links claim prices to replicating
portfolios and risk-neutral valuation.
\end{itemize}
\end{frame}

\begin{frame}{Q1(2): Black--Scholes Model}
\small
\textbf{Select all correct answers.} In the Black--Scholes model,
which statements are correct?
\begin{enumerate}
\item[A.] The underlying asset follows
$\dd S_t=\mu S_t\,\dd t+\sigma S_t\,\dd W_t$.
\item[B.] $W_t$ denotes a Wiener process.
\item[C.] The derivative value is written as $V_t=V(t,S_t)$.
\item[\wrong{D.}] \wrong{The volatility parameter is usually denoted by $r$.}
\item[E.] The risk-free bond grows as $B_t=B_0e^{rt}$.
\end{enumerate}
\end{frame}

\begin{frame}{Q1(2): Why?}
\small
\begin{itemize}
\item \textbf{A. True.} The stock is modeled by geometric Brownian motion.
\item \textbf{B. True.} $W_t$ supplies the Brownian randomness.
\item \textbf{C. True.} In a Markov model, the claim price depends on
current time and stock price.
\item \wrong{\textbf{D. False.} Volatility is $\sigma$; $r$ is the
risk-free interest rate.}
\item \textbf{E. True.} Constant continuously compounded interest
gives $B_t=B_0e^{rt}$.
\end{itemize}
\end{frame}

\begin{frame}{Q1(3): Binomial Tree}
\small
\textbf{Select all correct answers.} Which statements about the
binomial tree model are correct?
\begin{enumerate}
\item[A.] In one time step, the stock can move to $uS_i$ or $dS_i$.
\item[B.] The no-arbitrage condition is $d<e^{r\Delta t}<u$.
\item[C.] The risk-neutral probability is
$q=(e^{r\Delta t}-d)/(u-d)$.
\item[\wrong{D.}] \wrong{The real-world probability $p$ must be used
directly for derivative pricing.}
\item[E.] The derivative value can be computed by backward induction.
\end{enumerate}
\end{frame}

\begin{frame}{Q1(3): Why?}
\small
\begin{itemize}
\item \textbf{A. True.} The two branches are the defining stock moves.
\item \textbf{B. True.} The bond growth factor must lie strictly
between $d$ and $u$ to rule out arbitrage.
\item \textbf{C. True.} Solving
$e^{r\Delta t}=qu+(1-q)d$ gives this $q$.
\item \wrong{\textbf{D. False.} Pricing uses the risk-neutral
probability $q$; the real-world $p$ can differ.}
\item \textbf{E. True.} Start from terminal payoffs and discount
backward one step at a time.
\end{itemize}
\end{frame}

\begin{frame}{Q1(4): Trinomial Tree}
\small
\textbf{Select all correct answers.} Which statements about the
trinomial tree model are correct?
\begin{enumerate}
\item[A.] The stock price may move up, down, or remain unchanged.
\item[B.] In log-price form,
$x_{i+1}\in\{x_i+\Delta x,x_i,x_i-\Delta x\}$.
\item[C.] The probabilities are often denoted by $p_u,p_0,p_d$.
\item[D.] The probabilities satisfy $p_u+p_0+p_d=1$.
\item[\wrong{E.}] \wrong{The trinomial tree is stable for every choice
of $\Delta x$.}
\end{enumerate}
\end{frame}

\begin{frame}{Q1(4): Why?}
\small
\begin{itemize}
\item \textbf{A. True.} A trinomial step has three possible branches.
\item \textbf{B. True.} Equally spaced log-price nodes give these moves.
\item \textbf{C. True.} The subscripts indicate up, middle, and down.
\item \textbf{D. True.} The probabilities must sum to one; they must
also be nonnegative.
\item \wrong{\textbf{E. False.} Grid spacing affects the transition
weights. Some choices make a weight negative or the scheme unstable.}
\end{itemize}
\end{frame}

\begin{frame}{Q1(5): Finite Difference Stencil}
\small
\textbf{Select all correct answers.} Which finite difference
formulas and error claims are correct?
\begin{enumerate}
\item[A.] Forward difference:
$f'(x)\approx[f(x+h)-f(x)]/h$.
\item[B.] Centered first derivative:
$f'(x)\approx[f(x+h)-f(x-h)]/(2h)$.
\item[C.] Centered second derivative:
$f''(x)\approx[f(x-h)-2f(x)+f(x+h)]/h^2$.
\item[\wrong{D.}] \wrong{The forward difference has error $O(h^2)$.}
\item[E.] The centered first derivative has error $O(h^2)$.
\end{enumerate}
\end{frame}

\begin{frame}{Q1(5): Why?}
\small
\begin{itemize}
\item \textbf{A. True.} Taylor expansion gives
$[f(x+h)-f(x)]/h=f'(x)+O(h)$.
\item \textbf{B. True.} Subtract the Taylor expansions at
$x+h$ and $x-h$.
\item \textbf{C. True.} Add those expansions and subtract $2f(x)$.
\item \wrong{\textbf{D. False.} The forward difference is generally
first order, $O(h)$.}
\item \textbf{E. True.} Odd cancellation removes the $O(h)$ term,
leaving $O(h^2)$.
\end{itemize}
\end{frame}

\begin{frame}{Q1(6): Explicit Heat Scheme}
\small
For $u_\tau=u_{xx}$ with $\lambda=k/h^2$, \textbf{select all
correct answers} about the explicit scheme.
\begin{enumerate}
\item[A.] The next time level is computed directly from the current level.
\item[B.] One matrix form is
$u^{j+1}=(I-\lambda B)u^j+\lambda c^j$.
\item[\wrong{C.}] \wrong{A matrix equation must be solved at every
time step.}
\item[D.] The usual stability condition is $0<\lambda\leq1/2$.
\item[E.] Too large a $\lambda$ may make the numerical solution unstable.
\end{enumerate}
\end{frame}

\begin{frame}{Q1(6): Why?}
\small
\begin{itemize}
\item \textbf{A. True.} All right-hand values are already known.
\item \textbf{B. True.} With $B$ the tridiagonal negative-Laplacian
matrix, this is the explicit update; $c^j$ encodes boundaries.
\item \wrong{\textbf{C. False.} Explicit stepping is a matrix-vector
multiplication, not a linear solve.}
\item \textbf{D. True.} For the standard one-dimensional heat
equation, $\lambda\leq1/2$ controls the highest-frequency mode.
\item \textbf{E. True.} Beyond the stable range, errors can grow
rapidly during time stepping.
\end{itemize}
\end{frame}

\begin{frame}{Q1(7): Implicit Heat Scheme}
\small
\textbf{Select all correct answers.} Which statements about the
implicit finite difference scheme are correct?
\begin{enumerate}
\item[A.] It may require solving a linear system at each time step.
\item[B.] For the heat equation, one form is
$(I+\lambda B)u^j=u^{j-1}+\lambda c^j$.
\item[C.] Unknown new-level values appear in the equation.
\item[D.] $I+\lambda B$ is strictly diagonally dominant for $\lambda>0$.
\item[\wrong{E.}] \wrong{The implicit scheme is always identical
to the explicit scheme.}
\end{enumerate}
\end{frame}

\begin{frame}{Q1(7): Why?}
\small
\begin{itemize}
\item \textbf{A. True.} The new vector occurs inside a matrix system.
\item \textbf{B. True.} This is backward Euler for the heat
equation with boundary terms.
\item \textbf{C. True.} This is the reason the step is called implicit.
\item \textbf{D. True.} The diagonal is $1+2\lambda$, greater than
the sum $2\lambda$ of neighboring magnitudes.
\item \wrong{\textbf{E. False.} The two schemes evaluate spatial
differences at different time levels.}
\end{itemize}
\end{frame}

\begin{frame}{Q1(8): Consistency, Stability, Convergence}
\small
\textbf{Select all correct answers.}
\begin{enumerate}
\item[A.] Consistency means the scheme approximates the PDE locally
as the grid is refined.
\item[B.] Stability means numerical errors do not grow uncontrollably.
\item[C.] Convergence means the numerical solution approaches the exact
solution as $\Delta x,\Delta t\to0$.
\item[D.] The Lax equivalence theorem relates these notions for
suitable linear problems.
\item[\wrong{E.}] \wrong{A consistent scheme is always convergent,
even if it is unstable.}
\end{enumerate}
\end{frame}

\begin{frame}{Q1(8): Why?}
\small
\begin{itemize}
\item \textbf{A. True.} Substituting the exact solution leaves a
local residual that tends to zero.
\item \textbf{B. True.} Stability bounds the propagation of
perturbations over a fixed time interval.
\item \textbf{C. True.} Convergence concerns the global solution error.
\item \textbf{D. True.} For a well-posed linear problem, consistency
and stability are equivalent to convergence.
\item \wrong{\textbf{E. False.} A small local error can be amplified
without bound by an unstable scheme.}
\end{itemize}
\end{frame}

\begin{frame}{Q1(9): Lattices and Finite Differences}
\small
\textbf{Select all correct answers.}
\begin{enumerate}
\item[A.] Tree models can be interpreted as particular explicit
finite difference schemes.
\item[B.] A trinomial tree may arise from discretizing the
transformed Black--Scholes equation.
\item[C.] If $h^2=k\sigma^2$, then $p_0=0$, giving a binomial tree.
\item[D.] $p_u,p_0,p_d$ must be nonnegative to be probabilities.
\item[\wrong{E.}] \wrong{Lattice methods and finite difference methods
are completely unrelated.}
\end{enumerate}
\end{frame}

\begin{frame}{Q1(9): Why?}
\small
\begin{itemize}
\item \textbf{A. True.} Explicit updates combine values at neighboring
nodes, as trees do.
\item \textbf{B. True.} A centered space and forward time discretization
can produce three neighboring weights.
\item \textbf{C. True.} In that discretization,
$p_0=1-k\sigma^2/h^2=0$ when $h^2=k\sigma^2$.
\item \textbf{D. True.} Negative weights cannot represent
transition probabilities.
\item \wrong{\textbf{E. False.} The transition weights can be
exactly the finite difference update coefficients.}
\end{itemize}
\end{frame}

\begin{frame}{Q1(10): Fourier Stability Analysis}
\small
\textbf{Select all correct answers.}
\begin{enumerate}
\item[A.] Fourier analysis can study stability of constant-coefficient
finite difference schemes.
\item[B.] The amplification factor describes how a Fourier mode
evolves between time levels.
\item[C.] For the explicit heat scheme,
$g(\xi)=1-4\lambda\sin^2(\xi/2)$.
\item[D.] High-frequency errors must not grow uncontrollably.
\item[\wrong{E.}] \wrong{The explicit heat scheme is stable for all
$\lambda>0$.}
\end{enumerate}
\end{frame}

\begin{frame}{Q1(10): Why?}
\small
\begin{itemize}
\item \textbf{A. True.} Fourier modes diagonalize constant-coefficient
difference operators.
\item \textbf{B. True.} A mode is multiplied by $g(\xi)$ at each step.
\item \textbf{C. True.} Substitute $e^{\mathrm{i}i\xi}$ into the
explicit heat update.
\item \textbf{D. True.} The largest amplification often occurs near
the highest frequency.
\item \wrong{\textbf{E. False.} For the standard heat scheme,
$\lambda>1/2$ makes some modes grow.}
\end{itemize}
\end{frame}

\section{Question 2: CRR Tree}

\begin{frame}{Q2: Problem}
\small
Consider a European call option with strike price $K$ and expiry $T$.
The underlying stock follows geometric Brownian motion:
\[
\dd S_t=\mu S_t\,\dd t+\sigma S_t\,\dd W_t.
\]
Let the risk-neutral interest rate be $r$. Given
\[
S_0=K=100,\qquad r=0.05,\qquad \sigma=0.15,\qquad T=1,
\]
use a \textbf{two-step CRR binomial tree} to price the call.
\begin{enumerate}
\item[(a)] Compute the up factor $u$ and down factor $d$. \hfill (5 points)
\item[(b)] Compute the risk-neutral probability $p$. \hfill (5 points)
\item[(c)] Compute the value of the European call option. \hfill (5 points)
\end{enumerate}
\end{frame}

\begin{frame}{Q2(a): Up and Down Factors}
The inputs are
\[
S_0=K=100,\qquad r=0.05,\qquad \sigma=0.15,
\qquad T=1,\qquad N=2.
\]
The time step is $\Delta t=T/N=0.5$.

\begin{block}{CRR factors}
\[
u=e^{\sigma\sqrt{\Delta t}}\approx1.111895,
\qquad
d=e^{-\sigma\sqrt{\Delta t}}\approx0.899365,
\qquad ud=1.
\]
\end{block}

Check $d<e^{r\Delta t}<u$ before calculating the
risk-neutral probability.
\end{frame}

\begin{frame}{Q2(b): Risk-Neutral Probability}
The discounted stock price must have the same value before and after
one step:
\[
S=e^{-r\Delta t}\bigl(p\,uS+(1-p)dS\bigr).
\]
Solving for $p$ gives
\[
p=\frac{e^{r\Delta t}-d}{u-d}
=\frac{e^{0.025}-0.899365}{1.111895-0.899365}
\approx \boxed{0.592621}.
\]

This $p$ is denoted by $q$ in Tutorial 1.
The drift $\mu$ in the stock SDE does not enter the
no-arbitrage price of this option.
\end{frame}

\begin{frame}{Q2(c): Terminal Payoffs}
\begin{center}
\begin{tabular}{lcc}
\toprule
Path & Terminal stock price & Call payoff $(S_T-K)^+$\\
\midrule
$uu$ & $100u^2\approx123.6311$ & $23.6311$\\
$ud=du$ & $100ud=100$ & $0$\\
$dd$ & $100d^2\approx80.8858$ & $0$\\
\bottomrule
\end{tabular}
\end{center}
\vspace{0.5cm}
Only the $uu$ path finishes in the money.
The middle node recombines because $ud=du=1$.
\end{frame}

\begin{frame}{Q2(c): Backward Induction}
At the upper node after one step,
\[
C_u=e^{-r\Delta t}\bigl(pC_{uu}+(1-p)C_{ud}\bigr)
=e^{-0.025}p(23.6311).
\]
At the lower node, $C_d=0$. Discount one more step:
\[
\begin{aligned}
C_0
&=e^{-r\Delta t}\bigl(pC_u+(1-p)C_d\bigr)\\
&=e^{-rT}p^2(23.6311)
\approx\boxed{7.8945}.
\end{aligned}
\]
Keep unrounded $u,d,p$ until the final numerical answer.
\end{frame}

\section{Question 3: Explicit Scheme}

\begin{frame}{Q3: Problem (1/2)}
\small
Consider the initial boundary value problem
\[
\begin{cases}
u_t-2u_{xx}+u=0, &(x,t)\in(0,1)\times(0,\infty),\\
u(x,0)=\phi_0(x), &x\in(0,1),\\
u(0,t)=u(1,t)=0, &t>0.
\end{cases}
\]
Find a numerical solution on $[0,1]\times[0,T]$ for fixed $T>0$.
For $m,n\in\mathbb N$, set
\[
k=\frac{T}{m},\qquad h=\frac{1}{n+1},\qquad
x_i=ih,\qquad t_j=jk.
\]
Approximate $u_t(x_i,t_j)$ by a forward difference and
$u_{xx}(x_i,t_j)$ by a centered difference.
\end{frame}

\begin{frame}{Q3: Problem (2/2)}
\small
Let $v_{i,j}$ approximate $u(x_i,t_j)$. Answer the following:
\begin{enumerate}
\item[(a)] Write the finite difference scheme in terms of $v_{i,j}$.
\item[(b)] Find its amplification factor. Show stability with respect to
$\norm{\cdot}_{2,\Delta x}$, stating the condition if stability is
conditional.
\item[(c)] \begin{enumerate}
\item[(i)] Show consistency with the PDE in $\norm{\cdot}_\infty$,
stating the assumptions on the exact solution.
\item[(ii)] Write the scheme in matrix form and determine whether
it converges in $\norm{\cdot}_\infty$.
\end{enumerate}
\end{enumerate}
\end{frame}

\begin{frame}{Q3(a): Derive the Update}
Set $\lambda=k/h^2$. The zero boundary values are
$v_{0,j}=v_{n+1,j}=0$.
Insert the two difference formulas into the PDE:
\[
\frac{v_{i,j+1}-v_{i,j}}{k}
-2\frac{v_{i-1,j}-2v_{i,j}+v_{i+1,j}}{h^2}
+v_{i,j}=0.
\]
Solving for the next time level gives
\[
\boxed{
v_{i,j+1}
=2\lambda v_{i-1,j}
+(1-k-4\lambda)v_{i,j}
+2\lambda v_{i+1,j}.}
\]
The three weights sum to $1-k$; the reaction term $+u$ in
the PDE accounts for the loss of $k$.
\end{frame}

\begin{frame}{Q3(b): Amplification Factor}
Test one Fourier mode $v_{i,j}=\widehat v_j e^{\mathrm{i}i\xi}$.
The shifts multiply it by $e^{-\mathrm{i}\xi}$ and
$e^{\mathrm{i}\xi}$. Hence
\[
\begin{aligned}
g(\xi)
&=2\lambda e^{-\mathrm{i}\xi}
 +(1-k-4\lambda)+2\lambda e^{\mathrm{i}\xi}\\
&=1-k-4\lambda(1-\cos\xi)\\
&=\boxed{1-k-8\lambda\sin^2(\xi/2)}.
\end{aligned}
\]

For the finite interval with zero Dirichlet boundaries, the sine
modes have $\xi_\ell=\ell\pi/(n+1)$.
\end{frame}

\begin{frame}{Q3(b): The von Neumann Condition}
For stability over a fixed time interval, seek a constant $C$
independent of $h,k$ and $\xi$ such that
\[
|g(\xi)|\leq 1+Ck
\quad\text{for all admissible mesh sizes.}
\]

First control the diffusion part. Since
$0\leq\sin^2(\xi/2)\leq1$,
\[
-1\leq1-8\lambda\sin^2(\xi/2)\leq1
\quad\text{if}\quad \lambda\leq\frac14.
\]
Then
\[
|g(\xi)|
\leq\left|1-8\lambda\sin^2(\xi/2)\right|+k
\leq1+k.
\]
Thus $C=1$ works when $k/h^2\leq1/4$.
\end{frame}

\begin{frame}{Q3(b): Why $1+k$ Is Enough}
If $m$ steps reach time $T=mk$, then
\[
\norm{\bm v^{m}}_{2,h}
\leq (1+k)^m\norm{\bm v^0}_{2,h}
\leq e^T\norm{\bm v^0}_{2,h}.
\]
The constant $e^T$ stays bounded as the mesh is refined.

\begin{block}{Conditional stability}
The condition $k/h^2\leq1/4$ restricts the admissible meshes.
The explicit scheme is not unconditionally stable.
\end{block}

The stronger requirement $|g(\xi)|\leq1$ for every step gives
$k+8\lambda\leq2$. It is sufficient but is not needed for
the fixed-time stability argument above.
\end{frame}

\begin{frame}{Q3(c)(i): Local Truncation Error}
Insert the exact solution sampled on the grid into the difference
equation. At an interior point,
\[
\begin{aligned}
\tau_{i,j}
={}&\frac{u(x_i,t_j+k)-u(x_i,t_j)}{k}\\
&-2\frac{u(x_i-h,t_j)-2u(x_i,t_j)+u(x_i+h,t_j)}{h^2}
 +u(x_i,t_j).
\end{aligned}
\]
Taylor expansion gives
\[
\frac{u(t+k)-u(t)}{k}=u_t+O(k),
\]
\[
\frac{u(x-h)-2u(x)+u(x+h)}{h^2}
=u_{xx}+O(h^2).
\]
Because $u_t-2u_{xx}+u=0$, the leading PDE terms cancel.
\end{frame}

\begin{frame}{Q3(c)(i): Consistency in $\norm{\cdot}_\infty$}
If $u_{tt}$ and $u_{xxxx}$ are bounded on
$[0,1]\times[0,T]$, Taylor's theorem gives
\[
|\tau_{i,j}|
\leq \frac{k}{2}\norm{u_{tt}}_\infty
 +\frac{h^2}{6}\norm{u_{xxxx}}_\infty.
\]
Consequently,
\[
\boxed{\norm{\bm\tau^j}_\infty=O(k+h^2)\longrightarrow0.}
\]

The scheme is first order in time and second order in space.
The smoothness assumption is needed to make this bound uniform
over all grid points.
\end{frame}

\begin{frame}{Q3(c)(ii): Matrix Form}
Collect the $n$ interior values in
$\bm v^j=(v_{1,j},\ldots,v_{n,j})^\top$. The zero boundary
values give
\[
\bm v^{j+1}=A\bm v^j,
\qquad
A=\begin{pmatrix}
1-k-4\lambda & 2\lambda & & \\
2\lambda & 1-k-4\lambda & \ddots & \\
 & \ddots & \ddots & 2\lambda\\
 & & 2\lambda & 1-k-4\lambda
\end{pmatrix}.
\]
The diagonal coefficient contains both diffusion and reaction.
The off-diagonal coefficient is $2\lambda$ on each side.
\end{frame}

\begin{frame}{Q3(c)(ii): A Simple $\ell^\infty$ Condition}
To bound $\norm{A}_\infty$, make every entry of $A$
nonnegative. The only condition to check is
\[
1-k-4\lambda\geq0
\quad\Longleftrightarrow\quad
\boxed{k\leq\frac{h^2}{4+h^2}}.
\]
Each interior row sum is then $1-k$; boundary-adjacent rows
have a smaller sum. Hence
\[
\norm{A}_\infty\leq1-k<1,
\qquad
\norm{A^j}_\infty\leq1.
\]
This is a convenient \emph{sufficient} condition for stability
in the maximum norm.
\end{frame}

\begin{frame}{Q3(c)(ii): Convergence and a Weaker Bound}
The condition $\lambda\leq1/4$ from Q3(b) also gives
\[
\norm{A}_\infty
\leq |1-k-4\lambda|+4\lambda
\leq1+k,
\qquad
\norm{A^j}_\infty\leq e^{t_j}.
\]
If $\bm e^j$ is the grid error and the initial data are exact,
then $\bm e^0=0$ and
\[
\bm e^{j+1}=A\bm e^j-k\bm\tau^j.
\]
Therefore, for $t_j\leq T$,
\[
\norm{\bm e^j}_\infty
\leq T e^T\max_{\ell<j}\norm{\bm\tau^\ell}_\infty
=O(k+h^2)\longrightarrow0.
\]
Either maximum-norm bound proves convergence together with
the consistency estimate.
\end{frame}

\begin{frame}{Q3: Three Conditions to Keep Separate}
\begin{center}
\begin{tabular}{p{0.48\textwidth}p{0.40\textwidth}}
\toprule
Condition & What it supplies\\
\midrule
$\displaystyle k/h^2\leq1/4$ & A convenient von Neumann bound
$|g|\leq1+k$; also $\norm{A}_\infty\leq1+k$.\\[0.25cm]
$\displaystyle k\leq h^2/(4+h^2)$ & Nonnegative matrix entries and
$\norm{A}_\infty\leq1-k$.\\[0.25cm]
$\displaystyle k+8k/h^2\leq2$ & The stronger one-step bound
$|g|\leq1$ for the Fourier symbol.\\
\bottomrule
\end{tabular}
\end{center}
The second and third conditions are useful proof choices.
The exam does not require the sharpest stability boundary.
\end{frame}

\section{Question 4: Numerical Price}

\begin{frame}{Q4: Problem}
\small
The explicit finite difference demonstration, followed by linear
interpolation, gives a European call price at $t=0$ of
\$8.5889337. The Black--Scholes formula gives \$8.5916583,
a percentage error of about $0.0317\%$.

\medskip
Your boss wants to reduce this error. Give \textbf{two suggestions
with different origins}: one concerning the mathematical model and
one concerning the computational model or script.

\medskip
Explain each suggestion and any tradeoff. No mathematical proof or
programming is required. The demonstration script follows.
\end{frame}

\begin{frame}[fragile]{Q4: Demonstration 01 (1/2)}
\begin{lstlisting}
%Demonstration:
%Using Matlab to Implement Explicit Finite Difference Scheme
%to Price a European Call

clear;

%basic parameters
S0=100; K=S0; r=0.05; sigma=0.15; T=1;
n=25; % number of interior spatial points
m=50; % number of temporal steps (no. of iterations)

%computed constants
Sr=S0*exp((r-sigma^2/2)*T+6*sigma*sqrt(T));
h=Sr/(n+1); % spatial step size
k=T/m; %temporal step size

%define the vector b
b=zeros(n,1);
b(end)=0.5*(sigma^2*n^2+r*n)*Sr;
\end{lstlisting}
\end{frame}

\begin{frame}[fragile]{Q4: Demonstration 01 (2/2)}
\begin{lstlisting}
%Define the iteration matrix G
G=diag(-(sigma^2*(1:n)^2+r)); % G=A11
G=G+diag(0.5*(sigma^2*(1:(n-1))+r).*(1:(n-1)),1); % G=A11+A12
G=G+diag(0.5*(sigma^2*(2:n)-r).*(2:n),-1); % G=A11+A12+A21 (=A)
G=eye(n,n)+k*G; % G=I+kA

%Iteration
v=max(h*(1:n)'-K,0) % initialize v, v=v0

for j=1:m
    v=G*v+k*b; % v_j=G*v_(j-1)+k*b
end
\end{lstlisting}
\end{frame}

\begin{frame}{Q4: Where Can the Error Come From?}
The demonstration reports
\[
C_{\mathrm{FD}}=8.5889337,
\qquad C_{\mathrm{BS}}=8.5916583,
\qquad\text{relative error}\approx0.0317\%.
\]
The script uses $n=25$ interior spatial points, $m=50$
time steps, and a truncated stock-price domain $[0,S_r]$.

\begin{block}{Two distinct sources}
\begin{itemize}
\item Discretization and interpolation: finite $h$ and $k$.
\item Domain truncation: an artificial condition at $S=S_r$.
\end{itemize}
\end{block}
The task asks for one suggestion from each source, with tradeoffs.
\end{frame}

\begin{frame}{Q4: Refine the Computational Grid}
Increase $n$ to reduce spatial error and increase $m$ to reduce
time error. Linear interpolation at $S_0$ can also improve as
the spatial grid becomes finer.

\begin{block}{Explicit-scheme tradeoff}
If $h$ is roughly halved, keep $k/h^2$ within the stability
range by reducing $k$ roughly fourfold. This increases the
number of time steps and the computational cost.
\end{block}

Grid refinement should be checked by comparing prices on
successively finer grids; a single close result may contain
cancelling errors.
\end{frame}

\begin{frame}{Q4: Improve the Far-Field Boundary}
The script uses a constant boundary value $V(S_r,\tau)=S_r$.
For a European call, a better far-field approximation is
\[
\boxed{V(S_r,\tau)\approx S_r-Ke^{-r\tau}},
\qquad \tau=T-t.
\]
At maturity, $\tau=0$, this becomes $S_r-K$.
As the time-to-maturity grows, the strike is discounted.

\begin{block}{Implementation tradeoff}
The boundary contribution to the explicit update must now be
recomputed at each time step. Moving $S_r$ farther away can
also reduce boundary influence, but may require more spatial
points to preserve resolution near $S_0$ and $K$.
\end{block}
\end{frame}


\begin{frame}
\vfill
\begin{center}{\LARGE Thank You! Questions?}\end{center}
\vfill
\end{frame}

\end{document}
